\chapter{Der Raum als fraktale Energieverteilung}

Der Begriff des ``fraktalen Raumes'' ist irreführend. Physikalisch fundamental ist nicht die fraktale Geometrie, sondern die \emph{fraktale Verteilung von Energie}. Der Raum ist die Projektion dieser Verteilung.

Die fraktale Dimension \(D\) manifestiert sich nicht in der Metrik, sondern in den Quelltermen der Feldgleichungen. Für eine radialsymmetrische Massenverteilung gilt:

\begin{equation}
\rho_{\text{eff}}(r) = \rho_0 \left( \frac{r}{l_0} \right)^{D-3} = \rho_0 \cdot r^{-0,296}.
\end{equation}

Die Energiedichte ist also nicht homogen, sondern folgt einem Potenzgesetz.

\subsection{Drei Energiedichten, eine Skalierung}

Die Theorie kennt drei physikalisch relevante Energiedichten, die alle dem gleichen Skalierungsgesetz gehorchen:

\begin{table}[h]
\centering
\small
\setlength{\tabcolsep}{4pt}
\begin{tabular}{|p{2.8cm}|p{3.0cm}|p{3.8cm}|}
\hline
\textbf{Energieform} & \textbf{Skalierung} & \textbf{Physikalische Rolle} \\
\hline
Vakuum (negativ) & 
\(r^{D-3} \;( \approx r^{-0,296})\) & 
Quelle der Materieentstehung \\
\hline
Materie (positiv) & 
\(r^{D-3} \;( \approx r^{-0,296})\) & 
Galaxienhalos, Sterne, Gas \\
\hline
\(\Omega\)-Feld (kinetisch) & 
\(r^{D-3} \;( \approx r^{-0,296})\) & 
Gravitationswirbel, flache Rotationskurven \\
\hline
\end{tabular}
\caption{Die drei Energiedichten der WDBT+ und ihre gemeinsame Skalierung.}
\end{table}

\subsection{Die negative Vakuumenergiedichte}

Die negative Vakuumenergiedichte ist die fundamentale Eigenschaft des fraktalen Raumes:

\begin{equation}
\varepsilon_{\text{vac}}(r) = -\frac{\hbar c}{l_0^4} \left( \frac{l_0}{r} \right)^{3-D}.
\end{equation}

Sie ist \emph{nicht} die Dunkle Energie der ART. Sie ist eine lokale, räumlich variierende topologische Spannung des fraktalen Raumes, die die kontinuierliche Materieentstehung antreibt.

Die ``Löcher'' der fraktalen Struktur sind Regionen negativer Energiedichte – nicht Leere.